%% P07 --- Tensors, shapes and index notation
%%
%% Stub, generated by tools/gen_stubs.py from tools/programs.json.
%% The brief below is the contract for this program: what it must argue, and
%% what it must buy the reader. Writing the program means deleting the
%% \programstub{} block and replacing it with the program -- after which this
%% file is never regenerated. If the brief turns out to be wrong once the
%% mathematics has been worked, change it in the manifest and record the
%% contradiction in CLAUDE.md under *Resolved questions*.

\program{Tensors, shapes and index notation}
\label{prog:P07}

\programstub{%
Argument: an operation on arrays is a statement about which axes line up,
and index notation says it exactly where a picture cannot. Payoff: read an
\code{einsum} string as the sum it denotes; state the broadcasting rules and
say which axes align for a given pair of shapes; distinguish reshape from
transpose from permute, and say which of them moves data; count the axes of
an attention tensor (batch, head, query, key) and say what a matrix multiply
over them actually contracts. This is the largest gap in every book in the
field: the audience manipulates rank-4 arrays daily, has only ever been
taught rank-2 notation, and the mismatch is where the shape errors come
from.

\medskip
\textbf{Planned length:} 55 frames.
}
